% R&R Other Project Information, Field 10: Facilities & Other Resources
% (NOFO N0001426SF004, p. 20), for the copper/DED alternative
% (onr-yip-technical-proposal-cu-ded.tex). Text of the submission version of
% the evening of 24 Oct 2026 (agent-staging drop 2026-10-24-onr-yip-oct24h-merge,
% ee473d2, attachments-text/field10-facilities-text.md), which has no rates,
% prices or quote numbers. The submission PDF, with photographs, is built in
% byu-finances. This copy leaves the photographs out (some come from other
% institutions' web pages) and keeps their captions as plain text. Red TODOs
% mark what still needs Sterling, Prof. Fang or Mines.
% Build: xelatex field10-facilities-cu-ded.tex
\documentclass[12pt,letterpaper]{article}
\usepackage[margin=1in]{geometry}
\usepackage{fontspec}
\IfFontExistsTF{Times New Roman}{\setmainfont{Times New Roman}}{\setmainfont{TeX Gyre Termes}}
\usepackage{microtype}
\usepackage{xcolor}
\usepackage{xurl}
\usepackage[colorlinks=true,urlcolor=blue!70!black]{hyperref}
\setlength{\parindent}{0pt}
\setlength{\parskip}{5pt}
\pagestyle{plain}
\newcommand{\todo}[1]{\textcolor{red}{[TODO Sterling: #1]}}

\setlength{\emergencystretch}{3em}
\hyphenation{re-POW-DER AMA-ZE-MET}
\begin{document}
\begin{center}\textbf{Facilities and Other Resources}\\
ONR FY2027 Young Investigator Program, NOFO N0001426SF004. PI: Sterling G.\ Baird, Brigham Young University (BYU)\\
{\small\color{gray}Text copy: the photographs are left out and their captions kept as plain text.}\end{center}

\textbf{PI laboratory, BYU Department of Mechanical Engineering.} A 700
square foot powder metallurgy laboratory in the Clyde Building (CB 154),
assigned to the PI, holds an AMAZEMET rePOWDER ultrasonic atomizer with an
induction crucible module (Fig.~1b). We use the atomizer to make alloy powder
from elemental and master-alloy feedstock weighed into machined cups, about
four compositions per week and about 0.5~kg per run. For this project the
cups are oxygen-free copper and the powders are copper alloys. They are
melted at up to 1,400\,\textdegree{}C (typically up to
1,300\,\textdegree{}C) in graphite crucibles and atomized onto coated
carbon-carbon plates. For melts that contain nickel or iron, the crucible is
sprayed with boron nitride and the nozzle is boron nitride. A powder doser
built in the laboratory, an open-source auger doser that weighs each dose on
an analytical balance (0.1~mg), dispenses the powdered minor additions into
each charge and logs every dose. Sieves (ASTM B214), a Hall flowmeter funnel
(ASTM B212 and B213), glove bags and dry powder storage support the powder
checks. The atomizer stands in a steel-panelled enclosure built inside the
room in 2026 (Fig.~1a), with its own exhaust duct, argon supply, industrial
dehumidifier, electrical and chilled-water drops, a strip-curtain door and
tacky mats at the exit, so powder stays in the enclosure. The room also has
a canopy hood and a fume hood. The services were installed with the
engineering college's industrial hygienist and laboratory supervisor, who
reviewed the safety data sheets, the oxygen-deficiency risk of the argon
supply and the shipping of metal powders through BYU Chemicals Management;
a Class D extinguisher, an ATEX-rated (explosion-proof) vacuum cleaner,
full-face respirators and antistatic brushes are in place. It connects
directly to the BYU Electron Microscopy Facility by a door with no public
access.

\textbf{Figure 1.} The PI laboratory (CB 154). (a) The atomizer enclosure on
3 September 2026, after delivery. (b) The installed rePOWDER atomizer with
its induction module open, 28 September 2026. (c) Aluminum feedstock rods
melting in the induction-heated crucible. Frames from the laboratory's own
videos.

\textbf{Powder safety practice.} The laboratory's handling rules were set
for aluminum alloy powders, the most hazardous case in the Bureau of Mines
explosibility series, in conversations through 2026 with the powder group of
Ames National Laboratory (Dr.~Iver Anderson and Dr.~Jordan Tiarks: isolating
the collection can with a full-bore valve and opening it under inert gas,
controlled passivation, run length, the Bureau of Mines data), with
Professor Tim Horn of North Carolina State University (element-specific
chamber deposits, cleaning and shipping classification), with the group of
Professor Tim Weihs at Johns Hopkins University, which runs the same
atomizer for reactive powders, with the manufacturer's installer, with
Valimet, an industrial aluminum powder producer, and with the Manufacturing
Futures Institute at Carnegie Mellon University (powder-room rules). Three
trained operators run the atomizer; powder is handled in glove bags and
stored dry in sealed containers. Copper alloy powders are far less reactive,
and the same controls apply. Before any lot leaves BYU, a representative
powder of each new alloy family is tested for minimum ignition energy,
explosion severity and the UN N.1 and N.5 transport tests at an accredited
laboratory (Fauske \& Associates; in the budget), which set its handling
rules and shipping class.

\textbf{BYU Electron Microscopy Facility.} Scanning electron microscopes
with energy-dispersive X-ray spectroscopy (EDS) and electron backscatter
diffraction (EBSD) (Thermo Fisher Apreo C, Helios 600 with a focused ion
beam, and Verios), and a Tecnai TF-20 transmission electron microscope
(Fig.~2a). The instruments are shared and have high availability.

\textbf{Other BYU facilities.} Metallographic preparation uses PX500
automatic and PX300 manual polishers (Fig.~2c), vibratory polishers, an
etching fume hood and an electropolisher, with a Keyence VHX-7000 digital
optical microscope for inspection. We measure composition and structure at
the BYU X-ray Diffraction Facility and the BYU Spectroscopic Facility
(Agilent 8900 ICP-MS/MS, Fig.~2b, and XPS). The BYU Prototyping Lab machines
charge cups and test specimens by CNC. Sub-size tensile tests (ASTM
E8/E8M) run on the Instron load frames of the materials testing laboratory
in the Department of Mechanical Engineering \todo{confirm}. We measure
particle size distributions by sieve analysis in the PI laboratory and by
image analysis of SEM micrographs.

\textbf{Figure 2.} Shared BYU facilities used by the project. (a) Thermo
Fisher (FEI) Apreo C SEM, BYU Electron Microscopy Facility. (b) Agilent 8900
ICP-MS/MS, BYU Spectroscopic Facility. (c) PX300 manual polisher used for
metallographic preparation. Photos (a) and (b) are from the facilities' web
pages; (c) is from a laboratory training video.

\textbf{Computing.} The BYU Office of Research Computing operates a shared
cluster of about 28,000 CPU cores and 367 GPUs on 616 nodes
(\href{https://rc.byu.edu}{rc.byu.edu}, 2026), shared among BYU researchers
at no direct cost. CALPHAD calculations use Thermo-Calc with its copper alloy
database \todo{confirm the BYU license includes it} and the open-source
pycalphad; heat-flow calculations and the search software run on the same
cluster.

\textbf{Directed energy deposition.} Single-stream, blown-powder laser DED is
not available at BYU. The project buys machine time at two facilities, and
the search treats them as two steps of different cost. Neither facility is a
subawardee. \emph{The small step} is the InssTek MX-Lab (Fig.~3a) in the
laboratory of Professor Z.~Zak Fang (Department of Metallurgical
Engineering, University of Utah, Salt Lake City, 45 miles from BYU), which
his group has offered to the project. It has a 500~W fiber laser, 150~mm of
travel on each axis, control of melt-pool height, and six powder feeders
that can be run separately. It makes pre-placed tracks from a few grams of
powder and blown-powder tracks, thin walls and small blocks from about
0.5~kg, on a day trip from BYU. Its use is budgeted as hourly instrument
time until the group sets a rate \todo{confirm with Fang: the rate, the
smallest powder charge per feeder, the argon chamber and any preheat, and
who runs it}. \emph{The large step} is the FormAlloy X2 blown-powder laser
DED system (Fig.~3b) at the Colorado School of Mines (ADAPT center and
Shared Instrumentation Facility, Golden, Colorado), bought on a purchase
order. A typical session of four 25~mm cubes takes about 3 machine hours and
2 labor hours. The X2 takes powder from 50 to 200~$\mu$m, needs about 2~kg
of each powder to fill its hoppers, catches about half of the powder fed in
the melt pool, and has deposited copper-nickel tracks. It builds the tensile
and corrosion blocks and the repair coupons on thick plate, which need more
laser power than the MX-Lab has \todo{the X2's laser}. The United States
representative of InssTek has also offered deposition within the InssTek
network, which lists nickel-aluminum bronze (UNS C95300) among its qualified
materials.

\textbf{Figure 3.} The two DED steps. (a) The InssTek MX-Lab (500~W, six
powder feeders) at the University of Utah, the small step. (b) The FormAlloy
X2 (1~kW laser, four powder feeders, 250 $\times$ 250 $\times$ 300~mm build
envelope) at the Colorado School of Mines, the large step. Images from the
University of Utah metallurgy 3D printing center and the Mines ADAPT
capabilities pages.

\textbf{Corrosion testing.} The electrochemical screen runs on the
liquid-handling gantry in the PI laboratory (CubXL, Ursa Laboratories: a
pipette, a vial capper and two cameras on a scripted gantry with about
390 $\times$ 230 $\times$ 120~mm of travel). The gantry positions small
O-ring cells on polished specimens, with a NuVant EZstat Pro potentiostat
lent to the laboratory by a BYU colleague, which has run on the gantry since
October 2026. Immersion tests in substitute ocean water (ASTM D1141) run in
capped vials on the deck of the gantry, which samples the solutions for
ICP-MS at the BYU Spectroscopic Facility. Certified chemistry of powder lots
and final deposits (the alloying elements and impurities by ICP-OES and
ICP-MS; oxygen, nitrogen and hydrogen by inert gas fusion; carbon by
combustion) is purchased from NSL Analytical Services (Cleveland, Ohio; ISO
17025 and Nadcap accredited), the laboratory Ames Laboratory uses for its
atomized powders.

\textbf{Gas atomization.} The year's discovered alloy is converted into a
gas-atomized lot at Ames National Laboratory (Ames, Iowa), whose powder group
under Dr.~Jordan A. Tiarks (an unfunded collaborator) runs gas atomizers with
induction and plasma melting at scales from 10~g to 100~kg. One run a year
is bought as a service, inclusive of materials, equipment time and labor;
the laboratory is not a subrecipient. The gas-atomized lot is deposited at
Mines beside the ultrasonic lot of the same composition, so the project
measures what changes between the two routes, and it is the form in which
the final candidates go to Navy engineers.

\textbf{Government-owned facilities.} None are required. No additional
facilities will be acquired at the applicant's expense.

\end{document}
